% ═══════════════════════════════════════════════════════════════════════════
% LEHRERHANDREICHUNG DS-09: Einstimmungsstunde
% Tag VOR der Beobachtung -- Aktivierung ohne neue Inhalte
% ═══════════════════════════════════════════════════════════════════════════

\documentclass[11pt, a4paper]{article}

\usepackage[utf8]{inputenc}
\usepackage[T1]{fontenc}
\usepackage[ngerman]{babel}
\usepackage{helvet}
\renewcommand{\familydefault}{\sfdefault}

\usepackage[margin=1.2cm, top=1cm, bottom=1cm]{geometry}
\usepackage[table]{xcolor}
\usepackage{array}
\usepackage{tabularx}
\usepackage{booktabs}
\usepackage{tikz}
\usepackage{tcolorbox}
\usepackage{enumitem}
\usepackage{amssymb}
\tcbuselibrary{skins}

% Farben
\definecolor{spanisch}{HTML}{c41e3a}
\definecolor{gruppeA}{HTML}{2a7a4b}
\definecolor{gruppeB}{HTML}{1565c0}
\definecolor{hellgrau}{HTML}{f0f0f0}
\definecolor{phase}{HTML}{666666}
\definecolor{wichtig}{HTML}{b35c00}
\definecolor{gamification}{HTML}{6a0dad}

\pagestyle{empty}
\setlist{nosep, leftmargin=1.5em}

\begin{document}

% ═══════════════════════════════════════════════════════════════════════════
% HEADER
% ═══════════════════════════════════════════════════════════════════════════

\begin{tikzpicture}[remember picture, overlay]
\fill[spanisch] ([yshift=-1cm]current page.north west) rectangle ([yshift=-2.2cm]current page.north east);
\node[white, font=\Large\bfseries, anchor=west] at ([yshift=-1.6cm, xshift=1.2cm]current page.north west) {DS 09 -- Einstimmung (Tag vor Beobachtung)};
\node[white, font=\normalsize, anchor=east] at ([yshift=-1.6cm, xshift=-1.2cm]current page.north east) {90 Min | Mo 02.02.2026};
\end{tikzpicture}

\vspace{1.2cm}

% ═══════════════════════════════════════════════════════════════════════════
% WICHTIG-BOX
% ═══════════════════════════════════════════════════════════════════════════

\begin{tcolorbox}[
    colback=wichtig!15,
    colframe=wichtig,
    boxrule=1.5pt,
    top=4pt, bottom=4pt, left=6pt, right=6pt
]
\textbf{KEINE neuen Inhalte!} Bekanntes in neuer Kombination. Gamification für Motivation. Ziel: Selbstvertrauen stärken!
\end{tcolorbox}

\vspace{0.2cm}

% ═══════════════════════════════════════════════════════════════════════════
% ÜBERSICHT
% ═══════════════════════════════════════════════════════════════════════════

\begin{minipage}[t]{0.48\textwidth}
\textbf{Wiederholungsinhalte (bekannt):}
\begin{itemize}
\item \textit{ser + Adjektiv} (Eigenschaften)
\item \textit{estar + Adjektiv} (Zustände)
\item \textit{tener + años/hermanos}
\item \textit{hay que + Infinitiv}
\item Personenbeschreibung
\end{itemize}
\end{minipage}
\hfill
\begin{minipage}[t]{0.48\textwidth}
\textbf{Lernziele:}
\begin{itemize}
\item Vokabeln unter Zeitdruck abrufen
\item Spontan auf Prompts reagieren
\item Mitschüler beschreiben (4+ Sätze)
\item Stärken/Übungsfelder benennen
\end{itemize}
\end{minipage}

\vspace{0.2cm}
\hrule
\vspace{0.2cm}

% ═══════════════════════════════════════════════════════════════════════════
% VERLAUFSPLAN
% ═══════════════════════════════════════════════════════════════════════════

\textbf{\large Verlaufsplan}

\vspace{0.2cm}

\footnotesize
\begin{tabularx}{\textwidth}{|c|c|X|c|c|}
\hline
\rowcolor{spanisch!20}
\textbf{Zeit} & \textbf{Phase} & \textbf{Aktivität} & \textbf{SF} & \textbf{Material} \\
\hline
\multicolumn{5}{|c|}{\cellcolor{gamification!20}\textbf{PHASE 1: REPASO RELÁMPAGO (15 Min)}} \\
\hline
0--2 & Einstieg & Begrüßung: ``¡Hola! Hoy: Repaso divertido.'' & PL & -- \\
\hline
2--5 & Vorbereitung & 2 Teams bilden (A+B gemischt), Quiz-Regeln & PL & Tafel \\
\hline
5--14 & \textcolor{gamification}{\textbf{Quiz-Duell}} & 10 Fragen, Buzzer (Hand heben), Punkte & Teams & MATERIAL-09-quiz \\
\hline
14--15 & Auswertung & Siegerteam verkünden, Lob & PL & -- \\
\hline
\multicolumn{5}{|c|}{\cellcolor{hellgrau}\textbf{PHASE 2: LERNPFAD (25 Min)}} \\
\hline
15--17 & Überleitung & Lernpfad-Zugang, Geräte verteilen & PL & Tablets/Laptops \\
\hline
17--40 & Einzelarbeit & Selbstständig am Lernpfad (A2 oder B1 Track) & EA & LP-09 (digital) \\
\hline
40--42 & Abschluss LP & ``Wie weit seid ihr?'' -- Fortschritt melden & PL & -- \\
\hline
\multicolumn{5}{|c|}{\cellcolor{gamification!20}\textbf{PHASE 3: SPEED-DATING REMIX (20 Min)}} \\
\hline
42--45 & Vorbereitung & 2 Stuhlreihen, Mystery-Karten verteilen (je 3) & PL & MATERIAL-09-mystery \\
\hline
45--60 & \textcolor{gamification}{\textbf{Speed-Dating}} & 2 Min/Runde, Gong, wechseln, 5--6 Runden & PA rot. & Timer, Karten \\
\hline
60--62 & Auswertung & ``Was war die überraschendste Antwort?'' & PL & -- \\
\hline
\multicolumn{5}{|c|}{\cellcolor{gamification!20}\textbf{PHASE 4: TANDEM-CHALLENGE (20 Min)}} \\
\hline
62--65 & Vorbereitung & Gemischte Tandems (A+B), Challenge erklären & PL & Notizzettel \\
\hline
65--72 & Interview & Partner interviewen (3 Fragen), notieren & Tandem & -- \\
\hline
72--80 & Beschreibung & Partner beschreiben, Feedback geben & Tandem & -- \\
\hline
80--82 & Punkte & Beste Beschreibungen = Bonuspunkte & PL & Punkteliste \\
\hline
\multicolumn{5}{|c|}{\cellcolor{hellgrau}\textbf{PHASE 5: COOL-DOWN (10 Min)}} \\
\hline
82--85 & Reflexion & Selbst-Checkliste ausfüllen & EA & MATERIAL-09-checkliste \\
\hline
85--88 & Ausblick & ``Morgen kommt jemand zu Besuch -- normaler Unterricht!'' & PL & -- \\
\hline
88--90 & Abschluss & Positives Feedback, ``¡Hasta mañana!'' & PL & -- \\
\hline
\end{tabularx}

\normalsize
\vspace{0.2cm}

% ═══════════════════════════════════════════════════════════════════════════
% MATERIAL & TANDEM
% ═══════════════════════════════════════════════════════════════════════════

\begin{minipage}[t]{0.48\textwidth}
\begin{tcolorbox}[
    colback=hellgrau,
    colframe=phase,
    title={\footnotesize\bfseries Material-Checkliste},
    fonttitle=\bfseries,
    boxrule=0.5pt,
    top=2pt, bottom=2pt, left=4pt, right=4pt
]
\footnotesize
\begin{itemize}
\item[$\square$] \textbf{MATERIAL-09-quiz} (10 Fragen)
\item[$\square$] \textbf{MATERIAL-09-mystery-karten} (20 Karten)
\item[$\square$] \textbf{MATERIAL-09-checkliste} (10 Kopien)
\item[$\square$] Timer (sichtbar für alle)
\item[$\square$] Tablets/Laptops für LP-09
\item[$\square$] Stühle für Speed-Dating (2 Reihen)
\end{itemize}
\end{tcolorbox}
\end{minipage}
\hfill
\begin{minipage}[t]{0.48\textwidth}
\begin{tcolorbox}[
    colback=gamification!10,
    colframe=gamification,
    title={\footnotesize\bfseries Tandem-Zuordnung (Phase 4)},
    fonttitle=\bfseries,
    boxrule=0.5pt,
    top=2pt, bottom=2pt, left=4pt, right=4pt
]
\footnotesize
\begin{tabular}{@{}ll@{}}
\textcolor{gruppeA}{E.H.} + \textcolor{gruppeB}{A.S.} & (TOP + unsicher) \\
\textcolor{gruppeA}{L.N.} + \textcolor{gruppeB}{H.P.} & (proaktiv + Nachzügler) \\
\textcolor{gruppeA}{C.P.} + \textcolor{gruppeB}{A.A.} & (ruhig + unsicher) \\
\textcolor{gruppeA}{V.S.} + \textcolor{gruppeA}{L.K.} & (fleißig + abgelenkt) \\
\end{tabular}

\vspace{0.1cm}
\textit{E.H. als Vorbild, L.N. motiviert H.P.}
\end{tcolorbox}
\end{minipage}

\vspace{0.2cm}

% ═══════════════════════════════════════════════════════════════════════════
% AUSBLICK-FORMULIERUNGEN
% ═══════════════════════════════════════════════════════════════════════════

\begin{tcolorbox}[
    colback=white,
    colframe=spanisch,
    title={\footnotesize\bfseries Ausblick-Formulierungen (Phase 5)},
    fonttitle=\bfseries,
    boxrule=0.5pt,
    top=2pt, bottom=2pt, left=4pt, right=4pt
]
\footnotesize
\begin{minipage}[t]{0.48\textwidth}
\textbf{RICHTIG (positiv, entlastend):}
\begin{itemize}
\item ``Morgen kommt jemand zu Besuch.''
\item ``Ihr habt heute gezeigt, dass ihr das könnt!''
\item ``Wir machen einfach weiter wie immer.''
\item ``Es ist kein Test -- normaler Unterricht.''
\end{itemize}
\end{minipage}
\hfill
\begin{minipage}[t]{0.48\textwidth}
\textbf{FALSCH (Nervosität erzeugend):}
\begin{itemize}
\item[\texttimes] ``Morgen ist die wichtige Beobachtung.''
\item[\texttimes] ``Bitte benehmt euch morgen gut.''
\item[\texttimes] ``Ich werde bewertet, also...''
\item[\texttimes] ``Zeigt, was ihr könnt!'' (Leistungsdruck)
\end{itemize}
\end{minipage}
\end{tcolorbox}

\vspace{0.2cm}

% ═══════════════════════════════════════════════════════════════════════════
% SCHWIERIGKEITEN
% ═══════════════════════════════════════════════════════════════════════════

\begin{tcolorbox}[
    colback=white,
    colframe=wichtig,
    title={\footnotesize\bfseries Erwartete Schwierigkeiten},
    fonttitle=\bfseries,
    boxrule=0.5pt,
    top=2pt, bottom=2pt, left=4pt, right=4pt
]
\footnotesize
\begin{tabularx}{\textwidth}{@{}p{4cm}X@{}}
\textbf{SuS nervös wegen Beobachtung} & Fokus auf Spaß, nicht Perfektion. ``Morgen ist ein ganz normaler Tag.'' \\
\textbf{Quiz zu schnell für schwächere} & Einfache Fragen (1, 4, 6) gezielt an schwächere SuS richten. \\
\textbf{Technische Probleme LP} & AB-09-A/B (analog) als Backup bereithalten. \\
\textbf{Zeitproblem} & Speed-Dating auf 4 statt 6 Runden kürzen. Cool-Down kann gekürzt werden. \\
\end{tabularx}
\end{tcolorbox}

\end{document}
